\section{\ref{sec:neurontypes}장. 뉴런의 유형}

이 장의 수치는 사람의 뇌에서 세포를 직접 센 자료가 있는 곳에서는 그것에 기댄다. “뉴런 하나, 신경전달물질 하나” 규칙은 세포 아틀라스와 예외까지 찾아낸 실험에, \nt{DOPA}의 역할은 스파이크 기록에, 나머지 브로드캐스트 채널의 역할은 가설에 기댄다.

{\footnotesize
\setlength{\tabcolsep}{3pt}\renewcommand{\arraystretch}{1.15}
\begin{longtable}{>{\raggedright}p{0.5cm}>{\raggedright}p{4.0cm}>{\raggedright}p{5.6cm}>{\raggedright}p{2.3cm}>{\raggedright\arraybackslash}p{1.7cm}}
\toprule
No. & 명제 & 실험과 측정한 것 & 출처 & 상태 \\
\midrule\endfirsthead
\toprule
No. & 명제 & 실험과 측정한 것 & 출처 & 상태 \\
\midrule\endhead
1 & 사람의 뇌에는 뉴런이 약 $8.6\cdot10^{10}$개 있고, \nt{GLUT} 뉴런의 거의 전부는 소뇌의 작은 뉴런이다(표~\ref{tab:types}) & 등방성 분획법: 조직을 녹여 세포핵의 균일한 현탁액으로 만들고, 뉴런 표지 NeuN을 가진 핵을 센다. 성인 남성의 뉴런은 $86.1\pm8.1$\,십억(billion) 개이다. 그중 대뇌 피질에 있는 것은 $19\,\%$뿐이고 대부분은 소뇌에 있다 & \cite{Azevedo2009} & 측정됨 \\
2 & 거의 모든 뉴런에는 주된 신경전달물질이 하나 있지만(명제~\ref{prop:dale}) 예외도 있다 & 생쥐 뇌의 전사체 아틀라스: 세포 약 $4\cdot10^6$개, 클러스터 $5322$개. 합성과 포장 유전자에 따라 클러스터에 신경전달물질을 배정했다. 예외는 직접 측정되었다. 절편에서 \nt{DOPA} 뉴런은 같은 소포 수송체를 통해 GABA도 내보내 선조체 뉴런을 억제한다. 일부 \nt{DOPA} 축삭에서는 글루탐산과 도파민이 같은 전선의 서로 다른 말단에서 나온다(전자현미경과 광유전학) & \cite{Strata1999,Yao2023,Tritsch2012,Zhang2015} & 측정됨 \\
3 & 가중치 행렬의 열 전체가 같은 부호이다~\eqref{eq:dalesign} & 명제~\ref{prop:dale}와, 글루탐산 수용체는 거의 모두 흥분성이고 GABA 수용체는 억제성이라는 사실에서 장에서 유도한 것이다. “한 뉴런의 모든 수신자가 같은 부호의 가중치를 받는다”를 일반 규칙으로 직접 검증한 연구는 없다. 반례는 \nt{DOPA} 뉴런의 GABA 공동 방출(2행)과 부호가 다른 수용체를 가진 수신자들(11행)이다 & 이 장의 유도 & 이론 \\
4 & 피질 뉴런의 4분의 3에서 5분의 4는 \nt{GLUT}, 5분의 1에서 4분의 1은 \nt{GABA}이다 & 원숭이 피질 $10$개 영역에서 피질 전체 두께에 걸친 폭 $50$\,\textmu m의 기둥을 GABA로 면역염색했다. $8$개 영역에서 GABA 뉴런은 전체 뉴런의 약 $25\,\%$, 일차 시각 피질에서는 $20\,\%$ 미만이다 & \cite{Hendry1987} & 측정됨 \\
5 & \nt{GLUT} 뉴런의 출력은 약 $10^4$개이다 & 부검 입체해석학: 젊은 남성의 신피질에 시냅스 $1.64\cdot10^{14}$개(전자현미경, 디섹터법)와 뉴런 약 $2.3\cdot10^{10}$개. 비를 구하면 뉴런당 시냅스 $\sim7\cdot10^3$개이다. 평균적으로 입력 수는 출력 수와 같다 & \cite{Tang2001synapses,Pakkenberg1997} & 측정됨 \\
6 & \nt{DOPA} 뉴런의 출력은 말단 $10^5$--$10^6$개로 가지를 친다. 이것은 표적을 덮는 범위에서 얻은 추정치이지 센 값이 아니다 & 쥐의 \nt{DOPA} 뉴런 하나하나를 바이러스로 막 표지하고 축삭을 완전히 재구성했다. 뉴런 하나가 선조체 부피의 $0.45$에서 $5.7\,\%$(평균 $2.7\,\%$)를 덮는다. 측정한 것은 덮는 범위이지 말단 수가 아니다. 말단 수는 덮는 범위에서 자릿수 수준으로 유도했고, 말단을 직접 센 자료는 없다. \nt{SERO}와 \nt{NORA}의 말단 수는 거친 추정치이다 & \cite{Matsuda2009} & 측정됨 \\
7 & 브로드캐스트 뉴런은 적다. \nt{DOPA} $\sim5\cdot10^5$(두 무리 중 큰 쪽, 하한), \nt{SERO} $\sim3\cdot10^5$(두 부분 집계의 합), \nt{HIST} $\sim6\cdot10^4$(표~\ref{tab:types}, 그림~\ref{fig:typepie}) & 사람에서의 입체해석학적 집계: 흑색질의 색소 뉴런 $550\,000$개(배쪽 피개 영역 제외); 등쪽 솔기핵의 뉴런 $235\,000$개(핵의 모든 뉴런)와 다리뇌 피개의 세로토닌 뉴런 $125\,000$개; 시상하부의 히스타민 뉴런 약 $64\,000$개. \nt{NORA}, \nt{GLYC}, \nt{ACET}, \nt{ADRE}는 직접 센 자료가 없어 표에는 거친 자릿수 추정치를 넣었다 & \cite{Pakkenberg1991,Baker1990,Baker1991,Airaksinen1991} & 측정됨 \\
8 & 틈 속 글루탐산은 약 1밀리몰까지 치솟았다가 $\sim1\ms$ 만에 줄어든다 & 시냅스 전달 중 빠르게 떨어지는 경쟁적 차단제가 NMDA 수용체에서 밀려나는 정도로 측정했다(해마 뉴런 배양). 정점 $1.1$\,mM, 감소 시간 상수 $1.2\ms$ & \cite{Clements1992} & 측정됨 \\
9 & 작동 중인 피질에서는 흥분 전류와 억제 전류가 거의 균형을 이루며, 뉴런은 문턱값 근처에 머문다 & 이론: 연결이 강한 네트워크는 스스로 균형 체제에 이른다. 실험: 흰담비 전전두 피질 in vivo에서 “업” 상태 동안 시냅스 전류의 역전 전위는 수백 밀리초 동안 약 $-37$\,mV로 유지되는데, 컨덕턴스는 $\sim21$\,nS만큼 변한다. 흥분과 억제가 함께 오르내린다 & \cite{vanVreeswijk1996,Haider2006} & 측정됨 \\
10 & \nt{DOPA} 뉴런은 보상 예측 오차~\eqref{eq:rpe}를 알린다(그림~\ref{fig:rpe}) & 원숭이 \nt{DOPA} 뉴런의 스파이크: 예고 없는 주스에 버스트; 학습 뒤에는 신호에 버스트, 예측된 주스에는 무응답; 약속한 주스를 주지 않으면 발화율이 움푹 꺼진다. 정량 실험에서 발화율은 현재 보상과 과거 보상의 가중 평균의 차에 비례하지만, 양의 오차에서만 그렇다. 식~\eqref{eq:rpe} 자체는 시간차 알고리즘이다 & \cite{Schultz1997,Bayer2005,SuttonBarto2018} & 측정됨 \\
11 & 부호와 속도는 수용체가 정한다. 이온성은 1밀리초 미만, 대사성은 훨씬 느리다(명제~\ref{prop:receptor}) & 해마 뉴런의 막 조각에 글루탐산을 짧게 가했다. 이온성 수용체를 통한 전류는 $0.2$--$0.6\ms$ 만에 ($20$에서 $80\,\%$까지) 오르고 $2$--$3\ms$ 만에 줄어든다. 피라미드 뉴런의 느린 억제 전위는 대사성 GABA 수용체의 선택적 차단제로 사라진다. 도파민 수용체에는 작용이 반대인 두 계열이 있다. 선조체에서 D1 수용체를 가진 뉴런은 시냅스 강화에, D2를 가진 뉴런은 시냅스 약화에 도파민이 필요하다 & \cite{Colquhoun1992,DutarNicoll1988,Beaulieu2011,Shen2008} & 측정됨 \\
12 & 브로드캐스트 채널은 전선 한 가닥이 아니라 적은 수의 선로로 된 다발이다(\ref{sec:buses}절) & 생쥐 등쪽 솔기핵 세로토닌 뉴런의 바이러스-유전적 표지: 편도체로 출력을 보내는 뉴런과 이마엽 피질로 보내는 뉴런은 핵 안의 다른 자리에 있고, 서로 보완하는 영역으로 가지를 치며, 불쾌한 자극에 반대로 응답한다. 총설: 청반(\nt{NORA}) 뉴런의 하위 무리들은 서로 다른 과정을 부호화한다 & \cite{Ren2018,Poe2020} & 측정됨 \\
13 & \nt{NORA}는 이득 $g$를 정한다 & 적응적 이득 가설; 카테콜아민이 전달 특성의 기울기를 높이는 네트워크 모델. “노르아드레날린과 함께 $g$가 커진다”를 네트워크 전체에서 직접 측정한 자료는 없다. 원숭이에서 동공 지름이 청반 뉴런의 스파이크를 따라간다는 것은 측정되었다 & \cite{AstonJones2005,ServanSchreiber1990,Joshi2016} & 가설 \\
14 & \nt{ACET}은 “쓰기/읽기”를 전환하고 학습률을 정한다 & 가설. 근거: 쥐 후각 피질 절편에서 아세틸콜린 작용제($100$\,\textmu M)는 내부의 “회상” 시냅스를 $60\,\%$ 넘게 약화하지만 입력 시냅스는 $15\,\%$ 미만으로 약화한다 & \cite{Hasselmo2006,HasselmoBower1992} & 가설 \\
15 & \nt{SERO}는 할인율 $\gamma$를 정한다; 세로토닌 뉴런은 초당 스파이크 몇 개로 고르게 작동하고 수면의 한 단계에서는 거의 잠잠하다 & “세로토닌 $=\gamma$” 가설. 가장 강력한 근거: 생쥐 등쪽 솔기핵 세로토닌 뉴런을 광유전학으로 활성화하면 지연된 보상을 기다리다 일찍 포기하는 횟수가 줄어든다. 발화율과 급속 안구 운동 수면 중의 침묵은 측정되었다(기록 총설) & \cite{Doya2002,Miyazaki2014,JacobsAzmitia1992} & 가설 \\
\bottomrule
\end{longtable}}
